\documentclass[a4paper,11pt]{ltjsarticle}
\usepackage{geometry}
\geometry{top=25mm,bottom=25mm,left=25mm,right=25mm}
\usepackage{listings}
\usepackage{booktabs}
\usepackage{fancyhdr}
\usepackage{hyperref}
\usepackage{xcolor}
\usepackage{amsmath}
\usepackage{longtable}
\usepackage{array}
\usepackage{multirow}

\lstset{
  language=Verilog,
  basicstyle=\ttfamily\small,
  keywordstyle=\color{blue}\bfseries,
  commentstyle=\color{green!50!black},
  stringstyle=\color{red},
  numbers=left, numberstyle=\tiny\color{gray}, numbersep=5pt,
  frame=single, breaklines=true, tabsize=4, showstringspaces=false, captionpos=b
}

\pagestyle{fancy}
\fancyhf{}
\fancyhead[L]{計算機科学実験及演習3 ハードウェア — アーキテクチャ拡張仕様書}
\fancyfoot[C]{\thepage}
\renewcommand{\headrulewidth}{0.4pt}

\begin{document}

\begin{titlepage}
\centering
\vspace*{3cm}
{\LARGE 計算機科学実験及演習3 ハードウェア}\\[1cm]
{\Huge\bfseries アーキテクチャ拡張仕様書}\\[0.5cm]
{\LARGE SIMPLE/B プロセッサの拡張設計}\\[3cm]
\begin{tabular}{rl}
学籍番号 & （学籍番号を記入） \\[0.5em]
氏名     & （氏名を記入） \\[0.5em]
提出日   & 2026年6月5日 \\
\end{tabular}
\vfill
京都大学 工学部情報学科 計算機科学コース
\end{titlepage}

\tableofcontents
\clearpage


% =====================================================================
\section{拡張の概要}

本書では、SIMPLE/B基本設計に対して行ったアーキテクチャ拡張の仕様を定義する。
拡張は以下の5領域にわたる。

\begin{table}[h]
\centering
\caption{拡張項目一覧}
\begin{tabular}{clcl}
\toprule
\# & 拡張項目 & 担当 & 主な効果 \\
\midrule
1 & 4サイクル実行方式 & A & 性能25\%向上 \\
2 & ISA拡張 (BAL/BR/ADDI) & A & サブルーチン呼出、即値演算 \\
3 & 同期RAM (Block RAM) 対応 & A & メモリ容量256語→4096語 \\
4 & ハードウェアタイマ & B & 正確な時間計測 \\
5 & I/Oアドレス空間拡張 & B & 複数デバイス接続 \\
6 & 8桁7セグ表示+モード切替 & B & 32ビット表示、デバッグ支援 \\
\bottomrule
\end{tabular}
\end{table}


% =====================================================================
\section{マイクロアーキテクチャ拡張: 4サイクル化}

\subsection{基本設計との比較}

SIMPLE/Bの基本設計は5フェーズ順次実行であるが、
命令フェッチ（p1）とレジスタ書き戻し（p5）のデータパスが独立していることに着目し、
これらを1フェーズに統合して4フェーズとした。

\begin{table}[h]
\centering
\caption{フェーズ構成の比較}
\begin{tabular}{ll|ll}
\toprule
\multicolumn{2}{c|}{基本設計 (5フェーズ)} & \multicolumn{2}{c}{拡張設計 (4フェーズ)} \\
\midrule
p1 & 命令フェッチ (IF)       & \multirow{2}{*}{PH1} & \multirow{2}{*}{IF + WB} \\
p5 & レジスタ書き戻し (WB)   & & \\
p2 & レジスタ読み出し (RR)   & PH2 & RR \\
p3 & 演算 (EX)               & PH3 & EX + ST書込み \\
p4 & メモリアクセス/IO (MA)  & PH4 & MA / IO / 分岐 \\
\bottomrule
\end{tabular}
\end{table}

\subsection{フェーズ動作の詳細}

\begin{description}
\item[PH1 (IF+WB)]
前命令の書き戻し（パイプラインレジスタ経由）とフェッチを同時実行。
PH4で提示されたアドレスの命令を取得し、\texttt{ir}に格納。\texttt{pc $\leftarrow$ pc + 1}。

\item[PH2 (RR)]
\texttt{ar $\leftarrow$ r[Rs/Ra]}、\texttt{br $\leftarrow$ r[Rd/Rb]}。
PH1でのレジスタ書き込みはクロックエッジで完了しており、非同期リードで最新値が読まれる。

\item[PH3 (EX+ST)]
ALUまたはシフタで演算し、\texttt{dr}に格納。条件コード更新。
ST命令では\texttt{mem\_we=1}としPH3$\to$PH4エッジで書込み。

\item[PH4 (MA/IO/Branch)]
LD: \texttt{mdr $\leftarrow$ mem[dr]}。
IN/OUT: I/O操作。分岐: \texttt{pc $\leftarrow$ dr}。
書き戻し情報をパイプラインレジスタに保存。
次命令アドレスをRAMに提示。
\end{description}

\subsection{パイプラインレジスタ}

PH1で前命令の書き戻しを行うため、以下のレジスタをPH4で設定する。

\begin{table}[h]
\centering
\caption{書き戻しパイプラインレジスタ}
\begin{tabular}{lll}
\toprule
信号名 & 幅 & 用途 \\
\midrule
\texttt{wb\_rf\_we}   & 1 bit  & 書き込みイネーブル \\
\texttt{wb\_rf\_addr} & 3 bit  & 書き込み先レジスタ番号 \\
\texttt{wb\_rf\_data} & 16 bit & 書き込みデータ \\
\texttt{wb\_use\_mdr} & 1 bit  & LD/IN命令時、MDRを使用 \\
\bottomrule
\end{tabular}
\end{table}

\subsection{データハザードの回避}

レジスタファイルの非同期リードにより、PH1のクロックエッジでの書き込みと
PH2での読み出しは自然に整合する。フォワーディングやストールは不要。

\subsection{同期RAMとの協調}

Block RAM（altsyncram, \texttt{OUTDATA\_REG\_A = "UNREGISTERED"}）は
クロックエッジでアドレスをラッチし、組み合わせでデータを出力する。
データが必要な1フェーズ前にアドレスを提示する。

\begin{table}[h]
\centering
\caption{メモリアドレスのタイミング}
\begin{tabular}{lll}
\toprule
タイミング & \texttt{mem\_addr} & 目的 \\
\midrule
PH3 & \texttt{alu\_result} & LD/ST用実効アドレス \\
PH4 & \texttt{next\_pc\_val} & 次命令フェッチ \\
他  & \texttt{pc} & デフォルト \\
\bottomrule
\end{tabular}
\end{table}


% =====================================================================
\section{ISA拡張}

基本命令セットに3命令を追加する。全てformat (c)を使用し、
既存の未使用op2コードに割り当てる。

\begin{table}[h]
\centering
\caption{op2コード割り当て}
\begin{tabular}{lll}
\toprule
op2 & 命令 & 状態 \\
\midrule
000 & LI  & 基本 \\
001 & BAL & \textbf{拡張} \\
010 & BR  & \textbf{拡張} \\
011 & ADDI & \textbf{拡張} \\
100 & B   & 基本 \\
101 & (空き) & 将来拡張用 \\
110 & (空き) & 将来拡張用 \\
111 & Bcc & 基本 \\
\bottomrule
\end{tabular}
\end{table}

\subsection{BAL — Branch And Link}

\begin{table}[h]
\centering
\caption{BAL命令仕様}
\begin{tabular}{ll}
\toprule
項目 & 内容 \\
\midrule
形式        & (c): \texttt{10|001|Rb(3)|d(8)} \\
ニーモニック & \texttt{BAL Rb, d} \\
動作        & \texttt{r[Rb] = PC; PC = PC + 1 + sign\_ext(d)} \\
フラグ      & 変化なし \\
\bottomrule
\end{tabular}
\end{table}

PCはフェッチ時に+1済みのため、\texttt{r[Rb]}に保存されるのは次の命令のアドレスである。
ALU入力Aに\texttt{pc}を選択し、\texttt{dr = pc + sign\_ext(d8)}で分岐先を計算する。
書き戻しデータには\texttt{pc}（戻りアドレス）を使用する。

\subsubsection{コスト・ベネフィット分析}

\begin{itemize}
\item \textbf{回路コスト}: ALU入力Aマルチプレクサに1分岐追加、
  書き戻しデータの選択に1条件追加。追加ゲート数は約20 LE。
\item \textbf{ベネフィット}: サブルーチン呼出が1命令で実現。
  基本設計ではサブルーチン呼出は不可能であったため、質的な機能向上。
  プログラムの構造化、コード再利用が可能になる。
\end{itemize}

\subsection{BR — Branch Register}

\begin{table}[h]
\centering
\caption{BR命令仕様}
\begin{tabular}{ll}
\toprule
項目 & 内容 \\
\midrule
形式        & (c): \texttt{10|010|Rb(3)|d(8)} \\
ニーモニック & \texttt{BR Rb} \\
動作        & \texttt{PC = r[Rb] + sign\_ext(d)} （d は通常0）\\
フラグ      & 変化なし \\
\bottomrule
\end{tabular}
\end{table}

ALU入力Aに\texttt{br}（= \texttt{r[Rb]}）を選択。
レジスタ書き込み不要（\texttt{need\_rf\_write}に含めない）。

\subsubsection{コスト・ベネフィット分析}

\begin{itemize}
\item \textbf{回路コスト}: 分岐判定条件に\texttt{is\_br}を追加するのみ。約5 LE。
\item \textbf{ベネフィット}: BALと組み合わせてサブルーチンの復帰を実現。
  また、間接ジャンプ（ジャンプテーブル）にも使用可能。
\end{itemize}

\subsection{ADDI — Add Immediate}

\begin{table}[h]
\centering
\caption{ADDI命令仕様}
\begin{tabular}{ll}
\toprule
項目 & 内容 \\
\midrule
形式        & (c): \texttt{10|011|Rb(3)|d(8)} \\
ニーモニック & \texttt{ADDI Rb, d} \\
動作        & \texttt{r[Rb] = r[Rb] + sign\_ext(d)} \\
フラグ      & S, Z, C, V を更新 \\
\bottomrule
\end{tabular}
\end{table}

ALU入力Aに\texttt{br}（= \texttt{r[Rb]}）、入力Bに\texttt{sign\_ext(d8)}を選択。
既存ALU（ADD演算）をそのまま使用。

\subsubsection{コスト・ベネフィット分析}

\begin{itemize}
\item \textbf{回路コスト}: ALU入力Aマルチプレクサに1分岐追加、
  フラグ更新条件に\texttt{is\_addi}追加。約15 LE。
\item \textbf{ベネフィット}: ループカウンタのインクリメント/デクリメントが
  1命令で完了（基本設計では LI + ADD の2命令が必要）。
  典型的なループで33\%の命令数削減。
  さらにADDIのフラグ更新により、直後のBccで分岐判定が可能。
\end{itemize}

\subsection{拡張命令の効果まとめ}

\begin{table}[h]
\centering
\caption{ISA拡張による命令数削減}
\begin{tabular}{lcc}
\toprule
処理 & 基本ISA & 拡張ISA \\
\midrule
\texttt{r0 = r0 + 1} & 2命令 & 1命令 (ADDI) \\
関数呼び出し+復帰 & 不可能 & 2命令 (BAL+BR) \\
ループカウンタ操作+比較 & 3命令 & 2命令 (ADDI+CMP) \\
カウンタ操作+条件分岐 & 3命令 & 2命令 (ADDI+BNE) \\
\bottomrule
\end{tabular}
\end{table}


% =====================================================================
\section{メモリアーキテクチャ拡張: 同期RAM}

\subsection{基本設計との比較}

\begin{table}[h]
\centering
\caption{メモリの比較}
\begin{tabular}{lcc}
\toprule
項目 & 基本設計 & 拡張設計 \\
\midrule
メモリ種別 & Distributed RAM (非同期) & Block RAM (同期) \\
容量       & 256語 (8ビットアドレス) & 4096語 (12ビットアドレス) \\
読出し     & 組み合わせ & クロックエッジでラッチ \\
FPGA資源   & Logic Elements消費 & 専用RAMブロック使用 \\
\bottomrule
\end{tabular}
\end{table}

\subsection{メモリインタフェース}

\begin{lstlisting}[caption={RAM インタフェース}]
module ram(
    input  [11:0] address,  // 12ビットアドレス (4096語)
    input         clock,
    input  [15:0] data,     // 書込みデータ
    input         wren,     // 書込みイネーブル
    output [15:0] q         // 読出しデータ
);
\end{lstlisting}

Quartus Prime の altsyncram megafunction を使用し、
\texttt{run.mif} によりプログラムを初期化する。


% =====================================================================
\section{周辺回路拡張: ハードウェアタイマ}

\subsection{ブロック図}

\begin{figure}[h]
\centering
\begin{verbatim}
    +-------+     +-------------+     +-----------+
    | CTRL  |---->| プリスケーラ |---->| メイン     |---> COUNT (16bit)
    | [15:8]|     | カウンタ(8b) |     | カウンタ   |
    +-------+     +-------------+     +-----------+
    | EN(0) |--------+                    |   ^
    +-------+        |                    v   |
    | RST(1)|--------|-------> LIMIT比較 ---+
    +-------+        |           |
                     |           v
                     |       overflow (1bit) ---> Read-to-Clear
                     v
                 タイマ動作/停止
\end{verbatim}
\caption{タイマモジュールのブロック図}
\end{figure}

\subsection{レジスタ仕様}

\begin{table}[h]
\centering
\caption{タイマレジスタの詳細}
\begin{tabular}{lp{3cm}p{7cm}}
\toprule
レジスタ & ビット構成 & 動作 \\
\midrule
COUNT & [15:0] & プリスケーラで分周されたクロックでインクリメント。LIMIT到達で0リセット \\
CTRL  & [0] EN & 1で動作、0で停止 \\
      & [1] RESET & 1書込みでCOUNTを即座に0クリア（ワンショット、自動クリア） \\
      & [7:2] & 予約 \\
      & [15:8] PRESCALER & 分周値。カウント周波数 $= 20\,\mathrm{MHz} / (\mathrm{PRESCALER} + 1)$ \\
LIMIT & [15:0] & カウント上限。COUNT $\geq$ LIMITでoverflow発生、COUNT$\leftarrow$0 \\
overflow & [0] & 上限到達フラグ。IN命令で読み出すと自動クリア \\
\bottomrule
\end{tabular}
\end{table}

\subsection{インタフェース}

\begin{lstlisting}[caption={タイマモジュールのインタフェース}]
module timer(
    input         clk,
    input         rst_n,
    input  [15:0] wdata,       // 書き込みデータ
    input         we_ctrl,     // CTRL レジスタ書込み
    input         we_limit,    // LIMIT レジスタ書込み
    input         rd_overflow, // overflow 読出し (クリア用)
    output [15:0] count,       // 現在のカウント値
    output reg    overflow     // カウント到達フラグ
);
\end{lstlisting}


% =====================================================================
\section{周辺回路拡張: I/Oアドレス空間}

\subsection{設計方針}

IN/OUT命令のdフィールド下位4ビット（\texttt{I[3:0]}）を
I/Oポート番号として利用する。これにより最大16デバイスのアドレッシングが可能となる。

\subsection{CPUへの変更}

以下の3つの出力ポートを\texttt{simple\_cpu.v}に追加した。

\begin{lstlisting}[caption={CPUに追加したI/O制御信号}]
output [3:0]  io_port,    // I/Oポート番号 (= d4 = I[3:0])
output        io_read,    // IN命令実行中 (PH4)
output        io_write    // OUT命令実行中 (PH4)

assign io_port  = d4;
assign io_read  = (phase == PH4) && running && is_in;
assign io_write = (phase == PH4) && running && is_out;
\end{lstlisting}

\subsection{I/Oポートマップ}

\begin{table}[h]
\centering
\caption{I/Oポートマップ（拡張仕様）}
\begin{tabular}{cll}
\toprule
ポート & IN (読み出し) & OUT (書き込み) \\
\midrule
0  & DIPスイッチ A (8bit)   & 7セグ出力 下位16bit \\
1  & DIPスイッチ B (8bit)   & 7セグ出力 上位16bit \\
2  & タイマ COUNT (16bit)   & タイマ CTRL \\
3  & タイマ overflow (1bit) & タイマ LIMIT \\
4  & プッシュスイッチ (4bit)& LED出力 (8bit) \\
5  & ロータリースイッチ (4bit) & (予約) \\
6--15 & (予約: 0を返す)    & (予約: 無視) \\
\bottomrule
\end{tabular}
\end{table}

\subsection{後方互換性}

ポート番号0をDIPスイッチ/7セグに割り当てているため、
ポート番号を省略した既存プログラム（\texttt{d[3:0]=0}）は変更なしで動作する。


% =====================================================================
\section{周辺回路拡張: 8桁7セグメント表示}

\subsection{ダイナミック点灯の拡張}

桁選択信号を2ビット（4桁）から3ビット（8桁）に拡張した。
分周比4kHz を維持し、8桁巡回のリフレッシュレートは500Hzとなる。

\subsection{表示データの選択}

OUTポート0（下位16ビット）とポート1（上位16ビット）を結合した32ビットデータ、
またはデバッグ情報をロータリースイッチで切り替えて表示する。

\begin{table}[h]
\centering
\caption{表示モード}
\begin{tabular}{clp{7cm}}
\toprule
モード & データソース & 表示形式 \\
\midrule
0 & \texttt{\{seg\_reg\_hi, seg\_reg\_lo\}} & OUT命令の出力値を8桁HEX表示 \\
1 & \texttt{\{0000, debug\_pc\}} & PC値を下位4桁に表示 \\
2 & \texttt{\{r3[7:0], r2[7:0], r1[7:0], r0[7:0]\}} & r0\textasciitilde r3の下位8ビットを各2桁で表示 \\
3 & \texttt{\{0000, timer\_count\}} & タイマCOUNT値を下位4桁に表示 \\
\bottomrule
\end{tabular}
\end{table}

\subsection{デバッグポート}

CPU内部信号を外部から観測するため、以下の出力を追加した。

\begin{lstlisting}[caption={デバッグポートの実装}]
output [15:0] debug_pc,
output [15:0] debug_r0, debug_r1, debug_r2, debug_r3,

assign debug_pc = pc;
assign debug_r0 = rf.regs[0];  // 階層参照でレジスタファイル内部を参照
assign debug_r1 = rf.regs[1];
assign debug_r2 = rf.regs[2];
assign debug_r3 = rf.regs[3];
\end{lstlisting}


% =====================================================================
\section{モジュール構成}

\begin{table}[h]
\centering
\caption{最終的なモジュール構成}
\begin{tabular}{llp{7cm}}
\toprule
ファイル & モジュール & 概要 \\
\midrule
\texttt{simple\_top.v} & \texttt{simple\_top} & トップモジュール。全モジュール統合、I/Oデコード、8桁7セグ \\
\texttt{simple\_cpu.v} & \texttt{simple\_cpu} & CPUコア。4サイクルFSM、拡張ISA、I/Oポート出力 \\
\texttt{alu.v}         & \texttt{alu}         & 16ビットALU (7演算+フラグ) \\
\texttt{shifter.v}     & \texttt{shifter}     & バレルシフタ (4種) \\
\texttt{regfile.v}     & \texttt{regfile}     & 8$\times$16ビットレジスタファイル \\
\texttt{ram.v}         & \texttt{ram}         & Block RAM (altsyncram, 4096語) \\
\texttt{seg7dec.v}     & \texttt{seg7dec}     & 7セグデコーダ \\
\texttt{timer.v}       & \texttt{timer}       & ハードウェアタイマ \\
\bottomrule
\end{tabular}
\end{table}

\begin{figure}[h]
\centering
\begin{verbatim}
                        simple_top
    +------------------------------------------------------+
    |                                                      |
    |  +-----------+    +--------+    +--------+          |
    |  | simple_cpu|<-->|  ram   |    | timer  |          |
    |  |           |    |(4096w) |    |        |          |
    |  | +-----+   |    +--------+    +--------+          |
    |  | | alu  |  |                                      |
    |  | +-----+   |    +--------+                        |
    |  | |shift |  |    |seg7dec |---> seg[7:0]           |
    |  | +-----+   |    +--------+                        |
    |  | |regf  |  |                  sel[7:0] <---+      |
    |  | +-----+   |                               |      |
    |  +-----------+                  I/O Decoder --+      |
    |       ^   |                      ^    |              |
    |       |   v                      |    v              |
    |    in_mux  cpu_out     sw,pushsw,rotsw  led         |
    +------------------------------------------------------+
\end{verbatim}
\caption{システム全体のブロック図}
\end{figure}


% =====================================================================
\section{命令エンコーディング全体図}

\begin{table}[h]
\centering
\caption{命令エンコーディングマップ}
\small
\begin{tabular}{cc|l}
\toprule
op1 & op2/op3 & 命令 \\
\midrule
00 & --- & LD Ra, d(Rb) \\
01 & --- & ST Ra, d(Rb) \\
\midrule
10 & 000 & LI Rb, d \\
10 & 001 & BAL Rb, d \textbf{(拡張)} \\
10 & 010 & BR Rb \textbf{(拡張)} \\
10 & 011 & ADDI Rb, d \textbf{(拡張)} \\
10 & 100 & B d \\
10 & 101 & (空き) \\
10 & 110 & (空き) \\
10 & 111 & Bcc d (BE/BLT/BLE/BNE) \\
\midrule
11 & 0000 & ADD Rd, Rs \\
11 & 0001 & SUB Rd, Rs \\
11 & 0010 & AND Rd, Rs \\
11 & 0011 & OR Rd, Rs \\
11 & 0100 & XOR Rd, Rs \\
11 & 0101 & CMP Rd, Rs \\
11 & 0110 & MOV Rd, Rs \\
11 & 0111 & (空き) \\
11 & 1000 & SLL Rd, d \\
11 & 1001 & SLR Rd, d \\
11 & 1010 & SRL Rd, d \\
11 & 1011 & SRA Rd, d \\
11 & 1100 & IN Rd \\
11 & 1101 & OUT Rs \\
11 & 1110 & (空き) \\
11 & 1111 & HLT \\
\bottomrule
\end{tabular}
\end{table}


\end{document}
